\section{Ausgangslage}

Die Bundesverwaltung verfolgt eine Hybrid-Multi-Cloud Strategie, welche die Nutzung unterschiedlicher Cloud-Anbieter sowie der eigenen Infrastrukturen kombiniert. Das Bundesamt für Informatik und Telekommunikation (BIT) nimmt dabei die Rolle eines Cloud Service Brokers (CSB) wahr: Es stellt die Anbindung an verschiedene Hyperscaler sicher und betreibt gleichzeitig eine eigene On-Premises-Infrastruktur.

Für den Betrieb von Cloud-Diensten nutzen die einzelnen Provider sogenannte Landingzones - vordefinierte Strukturen, die den Rahmen für die Nutzung der jeweiligen Plattform bilden. Die bisherigen Erfahrungen zeigen jedoch, dass sich diese Landingzones zwischen den Anbietern grundlegend unterscheiden und bisher keine übergreifende technische Lösung existiert, die eine einheitliche Nutzung und Steuerung über mehrere Plattformen hinweg ermöglicht.

Damit entsteht die Herausforderung, Anforderungen und zentrale Elemente einer Provider-agnostischen Landingzone zu identifizieren, die eine konsistente, sichere und Compliance-konforme Nutzung der verschiedenen Plattformen innerhalb der Bundesverwaltung sicherstellt.

Im Rahmen des IP5 Projektes \textit{25fs imvs29} wurde sich zudem bereits mit der plattformübergreifenden Governance auseinandergesetzt. Mit einem PoC wurde dort gezeigt, dass eine agnostische Implementation grundsätzlich möglich ist. Die Erkenntnisse daraus haben möglicherweise Implikationen auf Teilaspekte der Arbeit in dieser Bachelorthesis.